\documentclass[runningheads]{llncs}

\usepackage[T1]{fontenc}

\begin{document}

\title{Individual Report}
\author{Baruch Shteken}
\institute{Universität Potsdam}

\maketitle

% ==================================================
\section{My Contribution}
% ==================================================

I contributed to the initial ideas for the project and was responsible
for most of the secondary encoding, including its Python implementation.
Due to my previous industry experience, I also contributed best practices
for Python development and GitHub workflows. When parts of the code behaved
differently from what we expected, I proposed possible explanations and
approaches for resolving these issues.

I was also responsible for running the experiments and documenting their
results. For the final paper, I wrote the section describing the Python
implementation as well as the results section. I also contributed to the
general development and refinement of the project throughout its different
stages.

% ==================================================
\section{Personal Reflection}
% ==================================================

I particularly liked the freedom we had to choose which aspects of the
project to explore. We could either focus on developing more complex
Clingo code or extend the Python implementation. At the same time, this
freedom was also one of the most challenging aspects of the project. We
initially had the idea of investigating track malfunctions, but we were
not sure how to develop this idea into a meaningful evaluation.

By looking at previous projects and papers from past students, we came
up with the idea of evaluating the running time of the secondary encoding.
This turned out to be a successful direction.

Another challenge was finding an effective cadence for working on the
project. In the beginning, we were unsure how frequently we should meet.
During the lecture period, the weekly meetings provided a good opportunity
to stay aligned. After the lecture period ended, we found that more
frequent, daily quick stand-ups were useful for keeping the project moving.

If I were to start the project again, I would establish regular weekly
meetings from the beginning and increase their frequency towards the end
of the semester as the deadline approached.

% ==================================================
\section{Recommendation to Future Students}
% ==================================================

Based on my experience, I would recommend that future students spend
some time identifying what they genuinely want to explore before
committing to a specific direction. We had two ideas at the beginning
and eventually abandoned one in favour of the other. This worked well
for us, but such a change should be made consciously rather than simply
because there is too little time left in the semester.

I would also recommend having at least one meeting per week where each
team member briefly explains what they have completed and what they are
currently working on. This makes it easier to identify when someone is
stuck and provides an opportunity for the team to help each other. Towards
the end of the semester, I would increase the frequency to daily or every
second day, particularly when approaching the final implementation and
evaluation.

\end{document}
